  \documentclass[12pt,german,english]{article}
    \usepackage[T1]{fontenc}
    \usepackage[ansinew]{inputenc}
%    \usepackage{a4}
   \usepackage{babel}
 %   \usepackage{german}
\usepackage{fullpage}

\begin{document}
\title{\bf  Whatever happened to monetarism?}
\author{Nikolaus K.A. L�ufer\\
Professor em., University of Konstanz, Germany} \vspace{-1.5mm}
\date{23.8.2009}
\maketitle

%\abstract{}

\section{Introduction: The Monetarist position revisited}
Thirty years ago Karl Brunner published a paper in which he described his brand of the "The Monetarist Revolution" in Monetary Theory.
He described his position by means of four defining characteristics: the nature of the transmission mechannism, the relative stability of the private sector,
the dominance of monetary impulses, the approximate separation between aggregative and allocative forces, and the irrelevance of the
endogeneity of the money stock.

The separability of aggregative and allocative forces is that part of this position which has come under little or no attack.
As a consequence, we have observed a demise in the use of large scale econometrics. Economietricians who once used to be
 admired for their work on large scale econometrics
have fallen in contempt. Combined with dramatic advances in time series analysis, large scale econometric models
gave way to small scale econometric time series models. In the meantime the structural approach has been reintroduced by
the backdoor of the residual term in time series analysis.


This is not a paper to review the defeat of the monetarist position.
Rather, it is an attempt to provide further evidence against the
monetarist position as outlined by Brunner. Actually, we shall
present some compelling evidence against the separability
hypothesis.

The starting point of my critique is a reformulation of the  import
function in conventional macroeconomic analysis. The reasoning in
support ot the new specification of the import function has been
given elsewhere.

\section{A new import function}
In a recent paper I have argued that the conventional import function of open economy macroeconomic models should
be reformulated.
The conventional import function is
\begin{equation}
 Im_t = Im(Y).
\end{equation}
This function should be replaced by the function
\begin{equation}
 Im_n = const + \kappa C(Y-T) + \alpha A + \iota I + \epsilon Ex .
\end{equation}

Accordingly, the income equilibrium condition for an open economy
\begin{equation}
 Y = C(Y-T) + I + A + Ex - Im(Y).
\end{equation}
should be substituted by
\begin{equation}
 Y = C(Y-T) + (1-\iota)I + (1-\alpha)A + (1-\epsilon) Ex - Im(Y-T)
\end{equation}
with a marginal propensity to import
\begin{equation}
m = \kappa \frac{\partial C}{\partial (Y-T)}
\end{equation}



\section{A critical thought experiment}
Suppose we change the exogenous I, A and Ex in such way a that their
sum $I + A + Es$ remains constant. Then
\begin{equation}
 d(I + A) = dI + dA = 0
\end{equation}

As long as we stick to the traditional import function, this change
cannot affect equilibrium income.
\begin{equation}
 d Y = 0
\end{equation}
If, instead, we apply the new import function, the condition for the
equilibrium income to remain constant for given changes in I, A, and
Ex is different:
\begin{equation}
(1-\iota)dI + (1-\alpha)dA = 0.
\end{equation}
Provided the import content paramters $\iota$, and $\alpha$ are
identically equal and smaller than one
\begin{equation}
\iota = \alpha < 1,
\end{equation}
then the condition for a zero change of equilibrium income is
equivalent to the former one:
\begin{eqnarray}
(1-\iota)dI + (1-\alpha)dA & = &(1-\alpha)(dI +
dA) \\
 & = & (dI + dA)=0
\end{eqnarray}
As long as the parameters $\iota$ and $\alpha$ are different, a
change in A and I such that $d(I+A) = dI + dA = 0$ must imply a
change in equilibrium income.

The demonstration can easily be extended to include changes of $Ex$.

\section{Final result}
As long as at least two of the import content parameters $\iota$,
$\alpha$, and $\epsilon$ are different, then nonzero changes of
$dI$, $dA$, and $dEx$ such that $dI + dA + dEx = 0$ will in general
affect equilibrium income. This contradicts Brunners separability
hypothesis of aggregative and allocative effects in macroeconommic
analysis. It reinstates the case for larger econometric models.
\bibliographystyle{plain}
\begin{thebibliography}{}
Nikolaus K.A. L�ufer, Neglecting the import content of shifts in
final demand - a major flaw in current macroeconomic analysis of
open economies, manuscript, August 2009
\end{thebibliography}

\end{document}
